% methodology.tex -- BNR-2026-013

\section{Methodology}
\label{sec:methodology}

\subsection{Literature Gate}

Following the laboratory protocol, we conducted a comprehensive
literature review of 31~sources before any experimental work:
12~primary papers from top cryptography venues (CRYPTO,
EUROCRYPT, USENIX Security, IACR ePrint), 6~production system
analyses (Scroll, Axiom, zkSync Era, Polygon zkEVM, Taiko,
Linea), 9~benchmark studies, and 4~implementation references.
Published benchmarks were recorded prior to experimentation to
establish calibration baselines.

\subsection{Elimination Analysis}

We apply constraint-based elimination against the targets in
Table~\ref{tab:constraints} before benchmarking. Any system
failing a hard constraint is eliminated regardless of
performance on other metrics. This avoids the need for
composite scoring with arbitrary weights.

\subsection{Benchmark Design}

\subsubsection{Hardware and Software}

All benchmarks were executed on a commodity desktop running
Windows~11 with Rust~1.93.0 in release mode (optimization
level~3, thin LTO). Library versions:

\begin{itemize}
  \item \textbf{Groth16}: arkworks (ark-groth16~0.5,
    ark-bn254~0.5)
  \item \textbf{halo2-KZG}: PSE fork v0.3.0 with KZG on BN254
    (halo2\_proofs, halo2curves~0.6)
\end{itemize}

\subsubsection{Circuit Design}

Two circuit families exercise distinct operational patterns:

\textbf{Arithmetic chain} (simulates EVM ADD/MUL): computes
$z = ((x^2 + x)^2 + (x^2 + x))^2 + \cdots$ for chain lengths
$\ell \in \{10, 50, 100, 500\}$. In Groth16, each step requires
1~R1CS constraint (multiplication; addition is free in R1CS). In
halo2, each step uses 1~row via a custom multiply-add gate with
constraint $s \cdot (a^2 + a - b) = 0$.

\textbf{Hash chain} (simulates Poseidon S-box $x^5$): computes
$z = ((x^5 + c_1)^5 + c_2)^5 + \cdots$ for chain lengths
$\ell \in \{10, 50, 100\}$. In Groth16, each round requires
3~R1CS constraints ($x^2$, $x^4 = (x^2)^2$, $x^5 = x^4 \cdot x$).
In halo2, each round uses 1~row via a custom degree-5 gate with
constraint $s \cdot (a^5 + c - b) = 0$.

\subsubsection{Statistical Protocol}

Each configuration runs 30~iterations after 2~warmup runs.
We report mean values. The benchmark framework outputs JSON
results with per-configuration setup time, proving time,
verification time, and proof size.

\subsection{Metrics}

We measure the following standard ZK metrics, each justified by
peer-reviewed precedent:

\begin{itemize}
  \item \textbf{Proving time} (ms): Wall-clock time for proof
    generation, the dominant cost in ZK
    rollups~\cite{zk-bottlenecks}.
  \item \textbf{Verification time} (ms): Local verification
    latency, proxy for on-chain gas cost.
  \item \textbf{Proof size} (bytes): Compressed proof size,
    determining on-chain calldata cost.
  \item \textbf{Constraint/row count}: Circuit complexity measure,
    directly proportional to proving time.
  \item \textbf{Setup type}: Per-circuit (Groth16) vs universal
    (halo2-KZG), qualitative operational metric.
\end{itemize}

On-chain verification gas is reported from published production
measurements (Axiom~\cite{axiom}, Nebra~\cite{nebra-gas},
Orbiter~\cite{orbiter-gas}) rather than simulated, as production
data is more authoritative than local EVM gas estimation.
